\documentclass[10pt,a4paper]{amsart}
\usepackage[margin=19mm]{geometry}
\usepackage{amsmath,amssymb,mathtools}
\usepackage[scr=rsfs,cal=euler]{mathalfa}
\usepackage{xfrac}
\usepackage[unicode,hidelinks,bookmarks=false]{hyperref}
\newcommand{\ie}{i.e.\ }
\newcommand{\cf}{cf.\ }
\newcommand{\resp}{respectively\ }
\newcommand{\Subsection}{\subsection}
\providecommand{\nobreakdash}{\nobreak\hbox{-}}
\setlength{\emergencystretch}{2em}
\multlinegap0pt
\usepackage{amssymb}
\newtheorem{definition}{Definition}
\newtheorem{theorem}{Theorem}
\def\R{\mathbb{R}}
\title{A sharp stability criterion for Euler equations via sparseness}
\author{\'{O}scar Dom\'{\i}nguez, Mario Milman}
\date{Source: arXiv:2310.19659v1 (CC BY 4.0)}
\begin{document}
\maketitle

\noindent\emph{Source and licence.} A sharp stability criterion for Euler equations via sparseness. Version arXiv:2310.19659v1 (CC BY 4.0), distributed by arXiv under the Creative Commons Attribution 4.0 International licence, http://creativecommons.org/licenses/by/4.0/, as recorded in the arXiv licence metadata for this exact version and shown on the abstract page https://arxiv.org/abs/2310.19659v1. Full licence notice: \texttt{licenses/arxiv-2310.19659-CC-BY-4.0.txt}.\par\smallskip

\noindent\textit{Editorial scope.} Counted unit: the article's theorem ThmSparseRieszMaximalRn (source0.tex lines 1391--1398). The article's complete original proof of it is delivered with it (source0.tex lines 1400--1720, one \texttt{proof} environment of 321 lines). No mathematics is written, repaired or re-derived: the statement and the proof are the article's own lines, in the article's own notation, and no retained line is substituted or re-flowed.  Retained uncounted as the source-local context the proof needs: lines 445--457; lines 996--1000; lines 1010--1027; lines 1012--1015; lines 1374--1378.  Nothing was omitted from any retained span, so the package carries no omission record.  No retained line contains a citation, so the wrapper carries no bibliography and no bibliography entry.  No retained line contains a figure environment, an image-inclusion command or drawing code, so no image is delivered and the package declares no asset dependency.  Editorial formatting only, in addition: an AMS \texttt{amsart} preamble that restates the article's own declarations and macros used by the retained lines, namely the environments \texttt{definition}, \texttt{theorem} on separate counters, together with the standard packages those lines need.  The retained environments print in this document as Definition 1, Theorem 1 in order of appearance, whereas the article prints the same items with the numbering of its own full document.  The complete native source of the article as distributed in the arXiv e-print for this exact version is bundled unchanged as \texttt{sources/149287/source0.tex}.
\begin{definition}[Sparse cubes]
\label{DefnSparse} A  (countable) family $(Q_{i})_{i\in
I} \subset \mathcal{D(}Q_{0})$ is called \emph{$\eta
$-sparse\footnote{In this paper the parameter $\eta$ will not play a role,
so in what follows we shall let $\eta=1/2$.}}, $\eta\in(0,1)$, if for every
$Q_{i}$ there exists a measurable subset $E_{Q_{i}}$ such that:
\begin{enumerate}
\item[(i)] the sets $E_{Q_{i}}$ are pairwise disjoint,
\item[(ii)] $\eta |Q_{i}|\leq|E_{Q_{i}}|$.
\end{enumerate}
We let $S(Q_{0})$ be the collection of all sparse families of
dyadic cubes in $Q_{0}.$ Analogously, one can introduce $S(\R^n)$, the set formed by all sparse  families of dyadic cubes in $\R^n$. 
\end{definition}

\begin{equation}
M_{\lambda,\alpha,Q_{0}}f(x)=\sup_{\substack{Q\in\mathcal{D}(Q_{0})\\x\in
Q}} |Q|^{\frac{\lambda}{n}-1}  (1-(\log|Q|)_{-})^{\alpha}\int_{Q}|f(y)| \,dy,\qquad x\in Q_{0}%
.\label{maxlam}%
\end{equation}

\begin{theorem}
\label{ThmSparseRieszMaximal} Suppose that $p,q,\alpha$ satisfy
\begin{equation}
1\leq p\leq q<\infty\qquad\text{and}\qquad\alpha\in\mathbb{R}\quad(\alpha
\leq0\text{ if }p=q). \label{AssumptionMon}%
\end{equation}
Then 
\[
SR_{p,q}\log^{\alpha}(Q_{0})=M_{n(\frac{1}{p}-\frac{1}{q}),\alpha,Q_{0}}%
L^{q}(Q_{0}).
\]
More precisely,
\begin{equation*}
\Vert f\Vert_{SR_{p,q}\log^{\alpha}(Q_{0})}\approx\Vert M_{n(\frac{1}{p}%
-\frac{1}{q}),\alpha,Q_{0}}f\Vert_{L^{q}(Q_{0})},
\end{equation*}
where the hidden constants of equivalence are independent of $f$ and $Q_{0}$.
\end{theorem}

\begin{equation}
1\leq p\leq q<\infty\qquad\text{and}\qquad\alpha\in\mathbb{R}\quad(\alpha
\leq0\text{ if }p=q). %
\end{equation}

\begin{equation}
M_{\lambda,\alpha}f(x)=\sup_{\substack{Q\in\mathcal{D}(\mathbb{R}^{n})\\x\in
Q}}|Q|^{\frac{\lambda}{n}-1}\big(1-(\log|Q|)_{-}\big)^{\alpha}\int%
_{Q}|f(y)|\,dy,\qquad x\in\mathbb{R}^{n}, \label{maxlamRn}%
\end{equation}

\begin{theorem}
\label{ThmSparseRieszMaximalRn} Suppose that $p,q, \alpha$ satisfy \eqref{AssumptionMon}.  Then
\[
SR_{p,q}\log^{\alpha}(\mathbb{R}^{n})=M_{n(\frac{1}{p}-\frac{1}{q}),\alpha
}L^{q}(\mathbb{R}^{n}).
\]

\end{theorem}

\begin{proof}
Consider the sequence of (not dyadic) cubes
\[
Q_{k} := [-2^{k}, 2^{k}]^{n}, \qquad k \in\mathbb{N}.
\]
According to Theorem \ref{ThmSparseRieszMaximal}, with equivalence constants
independent of $f$ and $k$,
\begin{equation}
\label{FatouMax0}\|f \mathbf{1}_{Q_{k}}\|_{SR_{p, q} \log^{\alpha}(Q_{k})}
\approx\|M_{n (\frac{1}{p} -\frac{1}{q}), \alpha, Q_{k}} (f \mathbf{1}_{Q_{k}%
})\|_{L^{q}(Q_{k})}.
\end{equation}


We claim that, for every $k\in\mathbb{N}$ and $x\in Q_{k}$,
\begin{equation}
M_{n(\frac{1}{p}-\frac{1}{q}),\alpha,Q_{k}}(f\mathbf{1}_{Q_{k}})(x)\approx
M_{n(\frac{1}{p}-\frac{1}{q}),\alpha}(f\mathbf{1}_{Q_{k}})(x),
\label{CorSparSobProof1}%
\end{equation}
(cf. \eqref{maxlam} and \eqref{maxlamRn}). Indeed, the estimate $\lesssim$
follows from the simple observation that $\mathcal{D}(Q_{k})\backslash
\{Q_{k}\}\subset\mathcal{D}(\mathbb{R}^{n})$, and the fact that the first
(dyadic) generation of $Q_{k}$, say $\{Q_{k,l}^{1}:l=1,\ldots,2^{n}\}$, gives
a pairwise disjoint decomposition of $Q_{k}$. In particular,
\[
Q_{k}=\bigcup_{l=1}^{2^{n}}Q_{k,l}^{1}%
\]
and $|Q_{k,l}^{1}|=2^{-n}|Q_{k}|$. Hence, given any $x\in Q_{k}$,
\begin{align*}
|Q_{k}|^{\frac{1}{p}-\frac{1}{q}-1}\,\int_{Q_{k}}|f(y)|\,dy  &  =2^{n(\frac
{1}{p}-\frac{1}{q}-1)}\,\sum_{l=1}^{2^{n}}|Q_{k,l}^{1}|^{\frac{1}{p}-\frac
{1}{q}-1}\int_{Q_{k,l}^{1}}|f(y)|\,dy\\
&  \leq2^{n(\frac{1}{p}-\frac{1}{q})}\,M_{n(\frac{1}{p}-\frac{1}{q}),\alpha
}(f\mathbf{1}_{Q_{k}})(x).
\end{align*}
Accordingly
\begin{align*}
M_{n(\frac{1}{p}-\frac{1}{q}),\alpha,Q_{k}}(f\mathbf{1}_{Q_{k}})(x)  &
\leq\sup_{\substack{Q\in\mathcal{D}(Q_{k})\backslash\{Q_{k}\}\\x\in
Q}}\,|Q|^{\frac{1}{p}-\frac{1}{q}-1}(1-(\log|Q|)_{-})^{\alpha}\,\int_{Q\cap
Q_{k}}|f(y)|\,dy\\
&  \hspace{1cm}+|Q_{k}|^{\frac{1}{p}-\frac{1}{q}-1}\,\int_{Q_{k}}|f(y)|\,dy\\
&  \lesssim\sup_{\substack{Q\in\mathcal{D}(\mathbb{R}^{n})\\x\in Q}%
}\,|Q|^{\frac{1}{p}-\frac{1}{q}-1}(1-(\log|Q|)_{-})^{\alpha}\,\int%
_{Q}|f(y)|\,\mathbf{1}_{Q_{k}}(y)\,dy\\
&  \hspace{1cm}+M_{n(\frac{1}{p}-\frac{1}{q}),\alpha}(f\mathbf{1}_{Q_{k}%
})(x)\\
&  \approx M_{n(\frac{1}{p}-\frac{1}{q}),\alpha}(f\mathbf{1}_{Q_{k}})(x).
\end{align*}


Next we focus on the estimate $\gtrsim$ in \eqref{CorSparSobProof1}. Consider
$x\in Q_{k},$ and $Q\in\mathcal{D}(\mathbb{R}^{n})$ with $Q\ni x$ and moreover
$Q\not \subset Q_{k}$ (indeed, if $Q\subset Q_{k}$ then $Q\in\mathcal{D}%
(Q_{k})$). We cannot assert that $Q_{k}\subset Q$ (recall that $Q_{k}$ is not
dyadic), but what is certainly true is that there exists $Q_{k}^{1}%
\in\mathcal{D}(Q_{k}),$ in the first dyadic generation (i.e., $2\,\ell
(Q_{k}^{1})=\ell(Q_{k})$), such that $Q_{k}^{1}\subset Q$ (because
$\mathcal{D}(Q_{k})\backslash\{Q_{k}\}\subset\mathcal{D}(\mathbb{R}^{n})$ and
$x\in Q_{k}$). Note that, in particular, $|Q|\geq|Q_{k}^{1}|=\ell(Q_{k}%
^{1})^{n}=2^{-n}\ell(Q_{k})^{n}=2^{-n}|Q_{k}|=2^{kn}>1$. Since $\frac{1}%
{p}-\frac{1}{q}-1<0,\ $we have
\begin{align*}
|Q|^{\frac{1}{p}-\frac{1}{q}-1}(1-(\log|Q|)_{-})^{\alpha}  &  =|Q|^{\frac
{1}{p}-\frac{1}{q}-1}\leq|Q_{k}^{1}|^{\frac{1}{p}-\frac{1}{q}-1}%
=2^{n(1+\frac{1}{q}-\frac{1}{p})}\,|Q_{k}|^{\frac{1}{p}-\frac{1}{q}-1}\\
&  =2^{n(1+\frac{1}{q}-\frac{1}{p})}\,|Q_{k}|^{\frac{1}{p}-\frac{1}{q}%
-1}(1-(\log|Q_{k}|)_{-})^{\alpha},
\end{align*}
which yields
\begin{align*}
|Q|^{\frac{1}{p}-\frac{1}{q}-1}(1-(\log|Q|)_{-})^{\alpha}\int_{Q\cap Q_{k}%
}|f(y)|\,dy  &  \leq 2^{n(1+\frac{1}{q}-\frac{1}{p})}\,|Q_{k}|^{\frac{1}{p}%
-\frac{1}{q}-1}(1-(\log|Q_{k}|)_{-})^{\alpha}\int_{Q_{k}}|f(y)|\,dy\\
&  \leq2^{n(1+\frac{1}{q}-\frac{1}{p})}\,M_{n(\frac{1}{p}%
-\frac{1}{q}),\alpha,Q_{k}}(f\mathbf{1}_{Q_{k}})(x).
\end{align*}
Consequently,
\[
M_{n(\frac{1}{p}-\frac{1}{q}),\alpha}(f\mathbf{1}_{Q_{k}})(x)\leq
2^{n(1+\frac{1}{q}-\frac{1}{p})}\,M_{n(\frac{1}{p}-\frac{1}{q}),\alpha,Q_{k}%
}(f\mathbf{1}_{Q_{k}})(x).
\]
This completes the proof of \eqref{CorSparSobProof1}.

On the other hand, since
\begin{equation}
\label{CubesStruct}Q_{k} \subset Q_{k+1} \qquad\text{and} \qquad\bigcup_{k
\in\mathbb{N}} Q_{k} = \mathbb{R}^{n},
\end{equation}
we have
\begin{equation}
\label{FatouMax1}\lim_{k \to\infty} M_{n (\frac{1}{p} - \frac{1}{q}), \alpha}
(f \mathbf{1}_{Q_{k}}) (x) \mathbf{1}_{Q_{k}} (x) = M_{n (\frac{1}{p} -
\frac{1}{q}), \alpha} f (x), \qquad x \in\mathbb{R}^{n}.
\end{equation}
Indeed, given any fixed $x \in\mathbb{R}^{n}$, we have
\[
M_{n (\frac{1}{p} - \frac{1}{q}), \alpha} (f \mathbf{1}_{Q_{k}}) (x)
\mathbf{1}_{Q_{k}} (x) \leq M_{n (\frac{1}{p} - \frac{1}{q}), \alpha} f (x),
\qquad\text{for all} \qquad k \in\mathbb{N},
\]
and (cf. \eqref{CubesStruct})
\begin{equation}
\label{FatouMax2}M_{n (\frac{1}{p} - \frac{1}{q}), \alpha} (f \mathbf{1}%
_{Q_{k}}) (x) \mathbf{1}_{Q_{k}} (x) \leq M_{n (\frac{1}{p} - \frac{1}{q}),
\alpha} (f \mathbf{1}_{Q_{k+1}}) (x) \mathbf{1}_{Q_{k+1}} (x).
\end{equation}
By the monotone convergence theorem for sequences of real numbers, we derive
\begin{align*}
\lim_{k \to\infty} M_{n (\frac{1}{p} - \frac{1}{q}), \alpha} (f \mathbf{1}%
_{Q_{k}}) (x) \mathbf{1}_{Q_{k}} (x)  &  = \sup_{k \in\mathbb{N}} M_{n
(\frac{1}{p} - \frac{1}{q}), \alpha} (f \mathbf{1}_{Q_{k}}) (x)\\
&  = \sup_{\substack{Q\in\mathcal{D}(\mathbb{R}^{n})\\x\in Q}}
|Q|^{\frac{1}{p}-\frac{1}{q}-1} \big(1- (\log|Q|)_{-} \big)^{\alpha} \sup_{k
\in\mathbb{N}} \int_{Q \cap Q_{k}}|f(y)| \,dy\\
&  = M_{n (\frac{1}{p} - \frac{1}{q}), \alpha} f (x),
\end{align*}
where we have used \eqref{CubesStruct} in the last step.

It follows from \eqref{CorSparSobProof1} that
\[
\Vert M_{n(\frac{1}{p}-\frac{1}{q}),\alpha,Q_{k}}(f\mathbf{1}_{Q_{k}}%
)\Vert_{L^{q}(Q_{k})}\approx\Vert M_{n(\frac{1}{p}-\frac{1}{q}),\alpha
}(f\mathbf{1}_{Q_{k}})\mathbf{1}_{Q_{k}}\Vert_{L^{q}(\mathbb{R}^{n})},
\]
uniformly with respect to $k$, consequently applying the monotone convergence
theorem (cf. \eqref{FatouMax1} and \eqref{FatouMax2}):
\begin{equation}
\lim_{k\rightarrow\infty}\Vert M_{n(\frac{1}{p}-\frac{1}{q}),\alpha,Q_{k}%
}(f\mathbf{1}_{Q_{k}})\Vert_{L^{q}(Q_{k})}\approx\Vert M_{n(\frac{1}{p}%
-\frac{1}{q}),\alpha}f\Vert_{L^{q}(\mathbb{R}^{n})}. \label{FatouMax3}%
\end{equation}


Next we deal with the left-hand side of \eqref{FatouMax0}. We claim that
\begin{equation}
\label{FatouMax4}\|f \mathbf{1}_{Q_{k}}\|_{SR_{p, q} \log^{\alpha}(Q_{k})}
\approx\|f \mathbf{1}_{Q_{k}}\|_{SR_{p, q} \log^{\alpha}(\mathbb{R}^{n})}%
\end{equation}
uniformly with respect to $k$ and $f$.

The estimate $\lesssim$ can be obtained as follows. Given any $(Q_{i})_{i\in
I}\in S(Q_{k})$, there are two possible scenarios. (I) $Q_{i}\neq Q_{k}$ for
every $i\in I$. Then $(Q_{i})_{i\in I}\in S(\mathbb{R}^{n}),$ since
$\mathcal{D}(Q_{k})\backslash\{Q_{k}\}\subset\mathcal{D}(\mathbb{R}^{n})$.
Clearly, this implies $\Vert f\mathbf{1}_{Q_{k}}\Vert_{SR_{p,q}\log^{\alpha
}(Q_{k})}\leq\Vert f\mathbf{1}_{Q_{k}}\Vert_{SR_{p,q}\log^{\alpha}%
(\mathbb{R}^{n})}$. (II) Suppose now that there is $i_{0}\in I$ such that
$Q_{i_{0}}=Q_{k}$. In particular, $(Q_{i})_{i\in I\backslash\{i_{0}\}}\in
S(\mathbb{R}^{n})$ and $Q_{i}\subset Q_{k}$ for $i\in I\backslash\{i_{0}\}$.
Now, the first dyadic decomposition of $Q_{k}$ (i.e., $\{Q_{k,l}%
^{1}:l=1,\ldots,2^{n}\}$) is formed by pairwise disjoint cubes in
$\mathcal{D}(\mathbb{R}^{n})$ (so, in particular, $\{Q_{k,l}^{1}%
:l=1,\ldots,2^{n}\}\in S(\mathbb{R}^{n})$) with $|Q_{k,l}^{1}|=2^{-n}|Q_{k}|$.
Hence, in this case, we can split the sum related to the $SR_{p,q}\log
^{\alpha}(Q_{k})$-norm as
\begin{align*}
\sum_{i\in I}\bigg(\frac{(1-(\log|Q_{i}|)_{-})^{\alpha}}{|Q_{i}|^{\frac
{1}{p^{\prime}}}}\,\int_{Q_{i}}|f(y)|\,dy\bigg)^{q}  &  =\\
&  \hspace{-5cm}\sum_{\substack{i\in I\\i\neq i_{0}}}\bigg(\frac
{(1-(\log|Q_{i}|)_{-})^{\alpha}}{|Q_{i}|^{\frac{1}{p^{\prime}}}}\,\int_{Q_{i}%
}|f(y)|\,dy\bigg)^{q}+\bigg(\frac{1}{|Q_{k}|^{\frac{1}{p^{\prime}}}}%
\,\int_{Q_{k}}|f(y)|\,dy\bigg)^{q}\\
&  \hspace{-5cm}\leq\Vert f\mathbf{1}_{Q_{k}}\Vert_{SR_{p,q}\log^{\alpha
}(\mathbb{R}^{n})}^{q}+\frac{1}{2^{\frac{nq}{p^{\prime}}}}\,\bigg(\sum
_{l=1}^{2^{n}}\frac{1}{|Q_{k,l}^{1}|^{\frac{1}{p^{\prime}}}}\int_{Q_{k,l}^{1}%
}|f(y)|\,dy\bigg)^{q}\\
& \hspace{-5cm}\lesssim \Vert f\mathbf{1}_{Q_{k}}\Vert_{SR_{p,q}\log^{\alpha
}(\mathbb{R}^{n})}^{q}+ \sum_{l=1}^{2^{n}} \bigg(\frac{1}{|Q_{k,l}^{1}|^{\frac{1}{p^{\prime}}}}\int_{Q_{k,l}^{1}%
}|f(y)|\,dy\bigg)^q\\
&  \hspace{-5cm}\lesssim\Vert f\mathbf{1}_{Q_{k}}\Vert_{SR_{p,q}\log^{\alpha
}(\mathbb{R}^{n})}^{q}.
\end{align*}
Therefore, taking the supremum over all possible $(Q_{i})_{i\in I}\in
S(Q_{k})$, we achieve
\[
\Vert f\mathbf{1}_{Q_{k}}\Vert_{SR_{p,q}\log^{\alpha}(Q_{k})}\lesssim\Vert
f\mathbf{1}_{Q_{k}}\Vert_{SR_{p,q}\log^{\alpha}(\mathbb{R}^{n})},
\]
i.e., the estimate $\lesssim$ in \eqref{FatouMax4} is shown.

To deal with the converse estimate, for any $\mathcal{Q} = (Q_{i})_{i \in I}
\in S(\mathbb{R}^{n})$, we consider the index set
\[
I_{k} := \{i \in I : Q_{i} \subset Q_{k}\}.
\]
Therefore we can split
\begin{align}
\sum_{i \in I} \bigg( \frac{(1-(\log|Q_{i}|)_{-})^{\alpha}}{|Q_{i}|^{\frac
{1}{p^{\prime}}}} \, \int_{Q_{i}} |f(y)| \, \mathbf{1}_{Q_{k}}(y) \, dy
\bigg)^{q}  &  =\nonumber\\
&  \hspace{-6cm} \sum_{i \in I_{k}} \bigg( \frac{(1-(\log|Q_{i}|)_{-}%
)^{\alpha}}{|Q_{i}|^{\frac{1}{p^{\prime}}}} \, \int_{Q_{i}} |f(y)| \, dy
\bigg)^{q}\nonumber\\
&  \hspace{-5cm}+ \sum_{i \in I \backslash I_{k}} \bigg( \frac{(1-(\log
|Q_{i}|)_{-})^{\alpha}}{|Q_{i}|^{\frac{1}{p^{\prime}}}} \, \int_{Q_{i} \cap
Q_{k}} |f(y)| \, dy \bigg)^{q}\nonumber\\
&  \hspace{-6cm} =: R_{1} + R_{2}. \label{SparseSumDecomp1}%
\end{align}


Note that $(Q_{i})_{i \in I_{k}} \in S(Q_{k})$ (since $(Q_{i})_{i \in I_{k}}
\in S(\mathbb{R}^{n}) \cap\mathcal{D}(Q_{k})$). Accordingly
\begin{equation}
\label{SparseSumDecomp2}R_{1} \leq\|f \mathbf{1}_{Q_{k}}\|_{SR_{p, q}
\log^{\alpha}(Q_{k})}^{q}.
\end{equation}


Concerning $R_{2}$, we argue as follows. Let $i\in I\backslash I_{k}$, i.e.,
$Q_{i}\not \subset Q_{k}$. Assume further that $Q_{i}\cap Q_{k}\neq\emptyset$.
Note that $Q_{k}$ is not a dyadic cube in $\mathbb{R}^{n}$, but its first
dyadic generation $\{Q_{k,l}^{1}:l=1,\ldots,2^{n}\}$ is formed by dyadic cubes
in $\mathbb{R}^{n}$. Since $Q_{k}$ can be expressed as the disjoint union of
the cubes $Q_{k,l}^{1}$, we can assert that there exists a unique $l(i)$ such
that $Q_{i}\cap Q_{k,l(i)}^{1}\neq\emptyset$. By the structure of dyadic cubes
in $\mathbb{R}^{n}$, we have either $Q_{i}\subset Q_{k,l(i)}^{1}$ or
$Q_{k,l(i)}^{1}\subset Q_{i}$. The former is not possible; otherwise,
$Q_{i}\subset Q_{k}$ but $i\not \in I_{k}$. Hence $Q_{k,l(i)}^{1}\subset
Q_{i}$, therefore
\begin{equation}
\int_{Q_{i}\cap Q_{k}}|f(y)|\,dy=\int_{Q_{k,l(i)}^{1}}|f(y)|\,dy.
\label{SparseSumDecomp3}%
\end{equation}


For $l \in\{1, \ldots, 2^{n}\}$, we define
\[
\mathcal{Q}_{k, l} := \{Q_{i} : i \in I \backslash I_{k} \quad\text{and} \quad
Q^{1}_{k, l} \subset Q_{i}\}.
\]
The above argument leads to
\begin{equation}
\label{SparseSumDecomp4}(Q_{i})_{i \in I \backslash I_{k}} = \bigcup
_{l=1}^{2^{n}} \mathcal{Q}_{k, l}.
\end{equation}
Moreover, since the $Q_{i}$'s are dyadic cubes in $\mathbb{R}^{n}$, the
elements of $\mathcal{Q}_{k, l} = \{Q_{1}, Q_{2}, \ldots\}$ can be ordered in
such a way that
$
Q^{1}_{k, l} \subset Q_{1} \subset Q_{2} \subset\ldots
$
We cannot exclude the possibility that some of the cubes in $\mathcal{Q}_{k,
l}$ coincide, but the number of these cubes is uniformly bounded by the sparse
constant $\eta$ in Definition \ref{DefnSparse}. Therefore, by
\eqref{SparseSumDecomp3} and \eqref{SparseSumDecomp4},
\begin{align}
R_{2}  &  \leq\sum_{i \in I \backslash I_{k}} \bigg( \frac{1}{|Q_{i}%
|^{\frac{1}{p^{\prime}}}} \, \int_{Q^{1}_{k, l(i)}} |f(y)| \, dy
\bigg)^{q}\nonumber\\
&  = \sum_{l=1}^{2^{n}} \bigg(\int_{Q^{1}_{k, l}} |f(y)| \, dy \bigg)^{q}
\sum_{Q_{i} \in\mathcal{Q}_{k, l}} \frac{1}{|Q_{i}|^{\frac{q}{p^{\prime}}}%
}\nonumber\\
&  \lesssim\sum_{l=1}^{2^{n}} \bigg(\int_{Q^{1}_{k, l}} |f(y)| \, dy \bigg)^{q}
\sum_{j=k}^{\infty}2^{-j \frac{n q}{p^{\prime}}}\nonumber\\
&  \approx\sum_{l=1}^{2^{n}} \bigg(\int_{Q^{1}_{k, l}} |f(y)| \, dy \bigg)^{q}
2^{-k \frac{n q}{p^{\prime}}}\nonumber\\
&  = \sum_{l=1}^{2^{n}} \bigg(\frac{1}{|Q^{1}_{k, l}|^{\frac{1}{p^{\prime}}}}
\int_{Q^{1}_{k, l}} |f(y)| \, dy \bigg)^{q}\nonumber\\
&  \leq\|f \mathbf{1}_{Q_{k}}\|^{q}_{SR_{p, q} \log^{\alpha}(Q_{k})}.
\label{SparseSumDecomp5}%
\end{align}


Combining \eqref{SparseSumDecomp1}, \eqref{SparseSumDecomp2} and
\eqref{SparseSumDecomp5},
\[
\sum_{i\in I}\bigg(\frac{(1-(\log|Q_{i}|)_{-})^{\alpha}}{|Q_{i}|^{\frac
{1}{p^{\prime}}}}\,\int_{Q_{i}}|f(y)|\,\mathbf{1}_{Q_{k}}(y)\,dy\bigg)^{q}%
\lesssim\Vert f\mathbf{1}_{Q_{k}}\Vert_{SR_{p,q}\log^{\alpha}(Q_{k})}^{q},
\]
for all $(Q_{i})_{i\in I}\in S(\mathbb{R}^{n})$. In particular,
\[
\Vert f\mathbf{1}_{Q_{k}}\Vert_{SR_{p,q}\log^{\alpha}(\mathbb{R}^{n})}%
\lesssim\Vert f\mathbf{1}_{Q_{k}}\Vert_{SR_{p,q}\log^{\alpha}(Q_{k})},
\]
completing the proof of \eqref{FatouMax4}.

By the lattice property of sparse RMT spaces, we have (recall
\eqref{CubesStruct})
\[
\Vert f\mathbf{1}_{Q_{k}}\Vert_{SR_{p,q}\log^{\alpha}(\mathbb{R}^{n})}%
\leq\Vert f\mathbf{1}_{Q_{k+1}}\Vert_{SR_{p,q}\log^{\alpha}(\mathbb{R}^{n})}%
\]
and
\[
\Vert f\mathbf{1}_{Q_{k}}\Vert_{SR_{p,q}\log^{\alpha}(\mathbb{R}^{n})}%
\leq\Vert f\Vert_{SR_{p,q}\log^{\alpha}(\mathbb{R}^{n})},\qquad k\in
\mathbb{N}.
\]
Hence, applying the monotone convergence theorem and the dominated convergence
theorem, we obtain
\begin{align}
\lim_{k\rightarrow\infty}\Vert f\mathbf{1}_{Q_{k}}\Vert_{SR_{p,q}\log^{\alpha
}(\mathbb{R}^{n})}  &  =\sup_{k\in\mathbb{N}}\Vert f\mathbf{1}_{Q_{k}}%
\Vert_{SR_{p,q}\log^{\alpha}(\mathbb{R}^{n})}\nonumber\\
&  =\sup_{(Q_{i})_{i\in I}\in S(\mathbb{R}^{n})}\sup
_{k\in\mathbb{N}}\bigg\{\sum_{i\in I}\bigg[\frac{\big(1-(\log|Q_{i}%
|)_{-}\big)^{\alpha}}{|Q_{i}|^{\frac{1}{p^{\prime}}}}\,\int_{Q_{i}\cap Q_{k}%
}|f|\bigg]^{q}\bigg\}^{\frac{1}{q}}\nonumber\\
&  =\sup_{(Q_{i})_{i\in I}\in S(\mathbb{R}^{n})}\lim
_{k\rightarrow\infty}\bigg\{\sum_{i\in I}\bigg[\frac{\big(1-(\log|Q_{i}%
|)_{-}\big)^{\alpha}}{|Q_{i}|^{\frac{1}{p^{\prime}}}}\,\int_{Q_{i}\cap Q_{k}%
}|f|\bigg]^{q}\bigg\}^{\frac{1}{q}}\label{SparseSumDecomp6}\\
&  =\sup_{(Q_{i})_{i\in I}\in S(\mathbb{R}^{n})}\bigg\{\sum
_{i\in I}\bigg[\frac{\big(1-(\log|Q_{i}|)_{-}\big)^{\alpha}}{|Q_{i}|^{\frac
{1}{p^{\prime}}}}\,\int_{Q_{i}}|f|\bigg]^{q}\bigg\}^{\frac{1}{q}}\nonumber\\
&  =\Vert f\Vert_{SR_{p,q}\log^{\alpha}(\mathbb{R}^{n}%
)}.\nonumber
\end{align}
Finally, taking limits on both sides of \eqref{FatouMax0} as $k\rightarrow
\infty,$ and invoking \eqref{FatouMax3} and \eqref{SparseSumDecomp6}, we
achieve the desired estimate
\[
\Vert f\Vert_{SR_{p,q}\log^{\alpha}(\mathbb{R}^{n})}\approx\Vert M_{n(\frac
{1}{p}-\frac{1}{q}),\alpha}f\Vert_{L^{q}(\mathbb{R}^{n})}.
\]

\end{proof}

\end{document}
